%% kapa: exam=1 round=35 year=2024 period=2 form=A subject=appraisal_law subject_ko=감정평가관계법규 seq=3 number=3
%% source: Q-Net 감정평가사 기출문제 (cst00309) artlSeq=5213229 file=raw/1008/5213229/2교시 문제지 원본.pdf
%% conversion: kordoc 4.24.51 -> tools/gen1.py (auto); review=proofread-sonnet; flags=none
\begin{question}{3}
국토의 계획 및 이용에 관한 법령상 도시ㆍ군관리계획에 관한 설명으로 옳지 않은 것은?
\begin{choices}
  \item 국토교통부장관은 국가계획과 관련된 경우에는 직접 도시ㆍ군관리계획을 입안할 수 있다.
  \item 도시ㆍ군관리계획은 광역도시계획과 도시ㆍ군기본계획에 부합되어야 한다.
  \item 주민은 지구단위계획의 수립에 관한 사항에 대하여 도시ㆍ군관리계획의 입안을 제안할 수 있다.
  \item 도시ㆍ군관리계획의 입안을 제안받은 자는 제안된 도시ㆍ군관리계획의 입안 및 결정에 필요한 비용의 전부를 제안자에게 부담시킬 수는 없다.
  \item 주거지역에 도시ㆍ군관리계획을 입안하는 경우 토지적성평가를 실시하지 아니할 수 있다.
\end{choices}
\end{question}
\answer{4}
